We have presented a systematic application of sparse autoencoders to single-cell foundation models,
producing feature atlases for two architecturally distinct models: 82,525 features across 18 layers
of Geneformer V2-316M and 24,527 features across 12 layers of scGPT whole-human. Both atlases
reveal models that have organized biological knowledge into modular, interconnected internal
representations. Against that organization stands a specific negative result: the features do not
carry transcription-factor-to-target regulatory logic at a rate distinguishable from what the
underlying measurements support.

\paragraph{Superposition is real but is not invisibility.}
Comparing thousands of dictionary atoms with a handful of principal axes cannot establish that a
model's structure is inaccessible to linear decomposition, because the two sides of that comparison
differ in capacity by two orders of magnitude. Under a matched comparison the picture is sharper and
more interesting. SAE directions are not principal axes---even against the complete
1,152-dimensional basis, under 1\% of features align with any single axis above a cosine of 0.5---but
neither are they hidden from the principal subspace: the median direction places half of its norm
in the leading 150--250 components, well ahead of the random-direction reference, and truncated SVD
needs rank 250--403 to match the SAE's own reconstruction quality. Nor is biological annotation the
exclusive property of non-aligned directions: annotated with an identical protocol, the leading 50
principal axes yield at least as many enriched terms per direction as randomly drawn SAE features.
What the dictionary provides is not access to otherwise unreachable information but a
\emph{different parameterization} of it---many sparse, one-sided, individually interpretable
directions in place of a few dense, signed, polysemantic ones. That is a real and useful property,
and it is the one worth claiming.

\paragraph{What a feature count does and does not mean.}
Dictionary atoms trained at different layers occupy different vector spaces and cannot be summed
into a count of concepts coexisting in one representation. Within a single layer the dictionary is
4$\times$ overcomplete and resolves roughly 1,250 distinct gene-level programs from 4,608 atoms; in
aggregate the 18 dictionaries hold 82,525 atoms that resolve 18,711 distinct programs, with 97\% of
multi-atom programs recurring in two or more layers. The interesting quantity is not a compression
ratio but this pattern of re-derivation: each layer rebuilds a largely fresh dictionary, yet the
gene-level programs those dictionaries express recur throughout the stack.

\paragraph{Hierarchical biological abstraction across layers.}
The U-shaped annotation profile, the shift in module themes from molecular machinery to integrative
programs, and the near-complete representational turnover between adjacent layers together describe
hierarchical biological abstraction. This parallels findings in large language models, where early
layers handle token-level processing and later layers more abstract semantic computation, but here
the tokens are genes and the semantics are biological programs.

\paragraph{Feature-level structure versus component-level nulls.}
Ablating a single feature produces a large, reproducible and gene-specific change in the model's
output where ablation of whole attention heads or MLP layers produces null behavioural
effects~\cite{kendiukhov2025systematic}. Computation is therefore organized around directions in
the residual stream rather than around architectural components, which is why component-level
interventions missed it. The scope of that effect is narrow, however: it falls on the genes the
feature fires on and not on the remainder of the biological process the feature is annotated with,
which is indistinguishable from a size- and activation-matched random gene set. A feature is a
handle on its own support, not on the pathway its label names---a distinction that matters for
anyone intending to use SAE features as intervention targets.

\paragraph{The co-expression--regulation boundary, and where it actually lies.}
Our results extend to the residual stream a boundary that has now been located by several
independent probes of the same models. Attention-derived edge scores add no incremental predictive
value over gene-level features for perturbation-target prediction~\cite{kendiukhov2025systematic}.
Causal circuit tracing recovers a dense, biologically coherent computational graph whose edges do
not correspond to directed regulatory relationships~\cite{kendiukhov2026circuittracing}, and
exhaustive mapping of that graph finds massive redundancy rather than identifiable regulatory
pathways~\cite{kendiukhov2026exhaustive}. Direct causal intervention on gene representations in
scGPT does not recover regulatory targets~\cite{kendiukhov2026causalintervention}, and adversarial
validation of in-silico perturbation profiles shows the same for predicted perturbation
responses~\cite{kendiukhov2026insilico}. Analyses of the representation geometry reach the same
place from a different direction: the structure these models learn is organized by cell state and
pathway membership~\cite{kendiukhov2026topology,kendiukhov2026celltypeprograms}, and although
residual-stream geometry does carry an incremental gene-regulatory signal, that signal is fragile
and does not survive stricter cross-validation~\cite{kendiukhov2026grsgeometry}. Sparse
dictionaries are one more probe, applied to the one representation the others did not decompose,
and they find what the others found.

No finite set of probes can establish that regulatory information is absent from a model, because
the proposition that it resides in some representation not yet examined is not refutable by
examination. What can be done is to state which representations have been examined, with which
instruments, and what was found. Attention weights, individual components, the causal graph
connecting them, gene-embedding geometry, in-silico perturbation responses, and now a sparse
decomposition of the residual stream have each been examined; none yields transcription
factor-to-target structure above what the underlying measurements support. That is the sense in
which we take the finding to be a property of these models rather than of any one instrument.

The important qualification is what that boundary can be attributed to, and here the evidence is
two-sided in a way that the primary dataset alone could not reveal.

A specificity rate is uninterpretable without knowing how much of the regulatory relationship is
recoverable from the same measurements. On the CRISPRi resource the atlas was built from, that
ceiling is very low: the assay measures 6,546 genes and the median perturbed transcription factor
has a single curated target among them. Differential expression computed directly on the knockdown
recovers curated targets for a minority of factors, established network inference recovers no more,
and the sparse features recover none, in any of four cell lines. The correct reading of that panel
is not that the models lack regulatory structure but that the measurement cannot decide the
question.

On perturbation data that measure the transcriptome, the same features do recover targets. In a
CRISPR-activation screen with 33,694 genes measured and strong induction, a factor's predicted
target set is enriched for its own high-confidence regulon in 15\% of cases, against 1\% for
perturbations of genes that are not in any reference network --- a fifteenfold separation from a
matched control, and therefore specific rather than a generic perturbation signature. Against
low-confidence regulons of several hundred genes the separation collapses and the apparent
enrichment is not specific, which is why we report the two confidence tiers separately.

The claim this study can support is therefore narrower than the absence of regulatory logic, and
wider than a null result confined to one assay. Sparse features of these models carry a small
amount of genuinely transcription-factor-specific target information, detectable only where the
assay measures the targets and the perturbation is strong, and negligible under the conditions in
which such analyses are usually run. Fifteen percent of a thirteen-factor panel is a modest signal
by any standard; it is also not nothing, and an interpretability method that reports zero on a
6,546-gene panel is measuring the panel as much as the model.

\paragraph{Where the unit of analysis matters, and by how much.}
A perturbation is applied to a cell, but a per-position statistic treats every gene position in
that cell as a separate observation. Whether this matters depends on how strongly feature
activation is correlated within a cell, which is an empirical question we can answer rather than
assume. The intraclass correlation of SAE feature activation across positions of the same cell is
$2.98 \times 10^{-4}$, which for the observed 2,020 positions per cell gives a design
effect of 1.61, that is an effective sample size of 62\% of the nominal count. The effective sample size is therefore a modest fraction of the
nominal position count rather than smaller by three orders of magnitude: sparse feature activations
at different gene positions of one cell are close to independent, because the sparsity pattern
dominates the within-cell variance.

The correction is nonetheless the right one to make, and we make it: significance is assessed with
the cell as the unit of replication throughout, holding the effect-size definition fixed at the
position-level scale so that the cell-level and position-level results differ only in what is
counted as an independent observation. We recommend the cell-level formulation as the default for
any interpretability analysis that tests per-position statistics against a per-cell intervention,
while noting that its practical effect here is bounded and measurable.

\paragraph{Implications for foundation model training.}
Current scFM training objectives may inherently bias representations toward co-expression. Learning
causal regulatory relationships would require training signals that distinguish cause from
correlation, such as perturbation-prediction objectives incorporated during pre-training. Equally,
evaluating whether a model has learned such relationships requires assays that measure the target
genes: a genome-scale perturbation atlas with a restricted expression panel bounds what any
interpretability method can demonstrate, however good the method is.

\paragraph{SAEs as a general interpretability framework for biological models.}
Independent of the regulatory question, this work establishes SAEs as a productive framework for
scFM interpretability. Applying an identical pipeline to two architecturally distinct models, with
qualitatively consistent results, demonstrates that these tools transfer to any transformer-based
biological model.

\paragraph{What would settle it.}
The decisive experiment is now clear and is not the one we set out to run. Regulatory structure in
these representations should be assessed on perturbation atlases that measure the whole
transcriptome, against high-confidence regulons, with perturbations of non-regulatory genes as a
matched control, and with the recoverable ceiling computed on the same cells. Applied to a
6,546-gene panel, any interpretability method will report a null result whose cause cannot be
identified. We would encourage the field to report the evaluability of the assay alongside any
claim about what a model does or does not encode.

\paragraph{Limitations.}
Several bounds apply to what we report. The regulatory analysis is limited by the gene panel of the
perturbation assay, as quantified above; a perturbation resource measuring the full transcriptome
would raise the ceiling against which any method is judged. TRRUST and DoRothEA are themselves
incomplete and literature-biased; a ChIP-seq-derived catalog would give an orthogonal view we have
not pursued. Our SAEs use one point in a wide architecture space, and although a targeted
expansion-by-sparsity ablation and a multi-seed stability analysis both leave the conclusions
intact, the space is not exhausted. Causal patching tests single-feature ablation; combinatorial
effects remain unexplored. Our SAEs are trained on unperturbed activations, which may bias the
learned dictionary toward cell-state and co-expression programs; training on pooled perturbed and
control activations with perturbation identity as an auxiliary signal is the most direct way to test
this and is beyond the scope of the present analysis. Finally, donor and assay identity are not
recorded in the cached extraction metadata for the Tabula Sapiens arm, so the batch analysis rests
on tissue as a coarse proxy.
